\chapter{Évaluation du Module 1}
\label{chap:evalmod1}

Cette évaluation valide l'acquisition des compétences liées à la gouvernance, l'analyse de risque et la prise de décision. Elle se décompose en trois parties : QCM, questions de synthèse et mise en situation pratique.

% ------------------------------------------------------------------------------
\section{Partie A : QCM (25 Questions)}
% ------------------------------------------------------------------------------

\textbf{Consigne :} Cochez la ou les bonnes réponses.

\begin{enumerate}
    \item La gouvernance du SI consiste principalement à :
    \begin{itemize}
        \item[A.] Installer des pare-feux.
        \item[B.] Définir les objectifs stratégiques et piloter la performance.
        \item[C.] Réparer les pannes informatiques.
        \item[D.] Programmer des applications métiers.
    \end{itemize}

    \item Dans une matrice RACI, qui est la personne qui "rend des comptes" et valide le résultat ?
    \begin{itemize}
        \item[A.] Responsible.
        \item[B.] Accountable.
        \item[C.] Consulted.
        \item[D.] Informed.
    \end{itemize}

    \item Quelle norme internationale traite spécifiquement de la gestion des risques en sécurité de l'information ?
    \begin{itemize}
        \item[A.] ISO 9001.
        \item[B.] ISO 27005.
        \item[C.] ISO 22301.
        \item[D.] ISO 26000.
    \end{itemize}

    \item Un "Actif" dans une analyse de risque est :
    \begin{itemize}
        \item[A.] Un logiciel de sécurité.
        \item[B.] Une ressource qui a de la valeur pour l'organisation.
        \item[C.] Un piratage informatique.
        \item[D.] Une norme obligatoire.
    \end{itemize}

    \item La formule de la Criticité est :
    \begin{itemize}
        \item[A.] Probabilité + Gravité.
        \item[B.] Probabilité $\times$ Gravité.
        \item[C.] Coût $\times$ Temps.
        \item[D.] Menace + Vulnérabilité.
    \end{itemize}

    \item Le risque résiduel est :
    \begin{itemize}
        \item[A.] Le risque avant toute mesure de protection.
        \item[B.] Le risque qui reste après application des mesures.
        \item[C.] Un risque qui ne peut jamais arriver.
        \item[D.] Un risque transféré à une assurance.
    \end{itemize}

    \item La méthode d'analyse de risque développée par l'ANSSI (France) s'appelle :
    \begin{itemize}
        \item[A.] NIST SP 800-30.
        \item[B.] OCTAVE.
        \item[C.] EBIOS RM.
        \item[D.] MEHARI.
    \end{itemize}

    \item Décider de ne pas héberger un service sensible pour ne pas prendre de risque est une stratégie de :
    \begin{itemize}
        \item[A.] Réduction.
        \item[B.] Acceptation.
        \item[C.] Transfert.
        \item[D.] Évitement.
    \end{itemize}

    \item Le RSSI est l'acronyme de :
    \begin{itemize}
        \item[A.] Responsable de la Sécurité des Systèmes d'Information.
        \item[B.] Réseau de Sécurité des Systèmes Informatiques.
        \item[C.] Responsable des Systèmes et Services Informatiques.
        \item[D.] Responsable Sécurité et Sureté Industrielle.
    \end{itemize}

    \item Une vulnérabilité est :
    \begin{itemize}
        \item[A.] Un actif critique.
        \item[B.] Une faiblesse pouvant être exploitée.
        \item[C.] Une attaque informatique.
        \item[D.] Un pare-feu désactivé.
    \end{itemize}

    \item Le DPO (Data Protection Officer) est principalement chargé de veiller au respect de :
    \begin{itemize}
        \item[A.] La loi de finances.
        \item[B.] Le RGPD.
        \item[C.] La sécurité physique.
        \item[D.] La comptabilité.
    \end{itemize}

    \item Dans EBIOS RM, l'atelier qui permet d'identifier les actifs essentiels est l'atelier :
    \begin{itemize}
        \item[A.] 1 (Contexte).
        \item[B.] 2 (Événements redoutés).
        \item[C.] 3 (Scénarios).
        \item[D.] 5 (Traitement).
    \end{itemize}

    \item Le triangle du risque est composé de :
    \begin{itemize}
        \item[A.] Actif, Menace, Vulnérabilité.
        \item[B.] Feu, Eau, Terre.
        \item[C.] Probabilité, Gravité, Coût.
        \item[D.] Hardware, Software, Humanware.
    \end{itemize}

    \item Souscrire une assurance cyber-risque est une stratégie de :
    \begin{itemize}
        \item[A.] Réduction.
        \item[B.] Acceptation.
        \item[C.] Transfert.
        \item[D.] Évitement.
    \end{itemize}

    \item L'alignement stratégique signifie que :
    \begin{itemize}
        \item[A.] La sécurité soutient les objectifs métier.
        \item[B.] Le métier obéit à la sécurité.
        \item[C.] La sécurité n'a pas de budget.
        \item[D.] L'IT et le Métier sont en conflit.
    \end{itemize}

    \item Le RACI est un outil de :
    \begin{itemize}
        \item[A.] Hacking.
        \item[B.] Répartition des rôles.
        \item[C.] Cryptographie.
        \item[D.] Développement logiciel.
    \end{itemize}

    \item La gravité d'un risque mesure :
    \begin{itemize}
        \item[A.] La chance qu'il se produise.
        \item[B.] L'ampleur des dommages s'il se produit.
        \item[C.] Le coût de la protection.
        \item[D.] La durée de l'attaque.
    \end{itemize}

    \item Le "COMEX" est :
    \begin{itemize}
        \item[A.] Un logiciel de gestion de mot de passe.
        \item[B.] Le Comité Exécutif de l'entreprise.
        \item[C.] Une norme ISO.
        \item[D.] Un type de virus.
    \end{itemize}

    \item Le PAS (Plan d'Action Sécurité) contient :
    \begin{itemize}
        \item[A.] Des actions, des responsables et des échéances.
        \item[B.] Des codes sources.
        \item[C.] Des mots de passe.
        \item[D.] Des factures fournisseurs.
    \end{itemize}

    \item Une erreur humaine (ex: envoi d'un fichier au mauvais destinataire) est :
    \begin{itemize}
        \item[A.] Une menace environnementale.
        \item[B.] Une menace humaine intentionnelle.
        \item[C.] Une menace humaine involontaire.
        \item[D.] Une vulnérabilité physique.
    \end{itemize}

    \item L'ISO 27001 concerne :
    \begin{itemize}
        \item[A.] La gestion de projet.
        \item[B.] Les exigences pour un SMSI (Système de Management de la Sécurité de l'Information).
        \item[C.] La qualité des logiciels.
        \item[D.] La continuité environnementale.
    \end{itemize}

    \item Le ROSPI (Return On Security Investment) mesure :
    \begin{itemize}
        \item[A.] La rentabilité financière de l'investissement sécurité.
        \item[B.] Le nombre d'attaques bloquées.
        \item[C.] La vitesse du réseau.
        \item[D.] La satisfaction utilisateur.
    \end{itemize}

    \item Un événement redouté concerne toujours l'impact sur :
    \begin{itemize}
        \item[A.] Le support technique.
        \item[B.] L'actif essentiel (Métier).
        \item[C.] Le budget.
        \item[D.] Le fournisseur.
    \end{itemize}

    \item Accepter un risque est justifié si :
    \begin{itemize}
        \item[A.] Le coût du traitement dépasse le coût de l'impact.
        \item[B.] La direction refuse tout budget.
        \item[C.] Le risque est critique.
        \item[D.] Le RSSI est absent.
    \end{itemize}

    \item Le cycle PDCA signifie :
    \begin{itemize}
        \item[A.] Plan, Do, Check, Act.
        \item[B.] Protect, Detect, Correct, Analyse.
        \item[C.] Plan, Deploy, Control, Audit.
        \item[D.] Program, Develop, Compute, Apply.
    \end{itemize}
\end{enumerate}

% ------------------------------------------------------------------------------
\section{Partie B : Questions Ouvertes (10 Questions)}
% ------------------------------------------------------------------------------

\textbf{Consigne :} Répondez de manière synthétique (5-10 lignes par question).

\begin{enumerate}
    \item Expliquez la différence fondamentale entre "Gouvernance" et "Management" dans le contexte de la cybersécurité.
    \item Pourquoi le RSSI ne doit-il pas être rattaché hiérarchiquement à la DSI ? Quel est le risque de conflit d'intérêts ?
    \item Définissez les termes "Actif", "Menace" et "Vulnérabilité" et donnez un exemple concret pour chacun dans un contexte de banque en ligne.
    \item Quelle est la différence entre le "Risque Inhérent" et le "Risque Résiduel" ?
    \item Pourquoi la méthode EBIOS RM est-elle particulièrement adaptée aux entreprises françaises ?
    \item Expliquez le rôle du "Accountable" dans une matrice RACI. Pourquoi est-il unique pour chaque tâche ?
    \item Citez les quatre options de traitement du risque et donnez un exemple pour chacune dans le contexte d'un risque de vol d'ordinateur portable.
    \item Qu'est-ce que le ROSPI et pourquoi est-il utile pour convaincre le COMEX ?
    \item Décrivez les critères D.I.C.T (Disponibilité, Intégrité, Confidentialité, Traçabilité) avec un exemple pour chacun.
    \item À quoi sert un Plan d'Action Sécurité (PAS) ? Quels sont ses éléments constitutifs minimums ?
\end{enumerate}

% ------------------------------------------------------------------------------
\section{Partie C : Cas Pratique - Mise en Condition (TP)}
% ------------------------------------------------------------------------------

\subsection{Scénario : La Menace Interne}

\begin{boxlizbiz}
\textbf{Contexte :}
Nous sommes 3 mois après votre première mission. Le projet de segmentation réseau est validé. Cependant, ce matin, la DSI de LiZbiZ vous appelle en urgence.

\textbf{Situation :} Un ancien développeur, parti en mauvais termes la semaine dernière, a tenté d'accéder au serveur de code source hier soir. Son compte n'avait pas été désactivé. L'alerte a été levée par le système de détection (SIEM), mais l'intrusion a échoué de justesse car le mot de passe du serveur était complexe.

Le Directeur Général est inquiet : "Si ce développeur avait réussi, il aurait pu voler notre code source ou insérer une porte dérobée !"
\end{boxlizbiz}

\subsection{Travaux à Réaliser}

Vous avez 1h30 pour produire un rapport complet (Format Libre : PDF, PowerPoint).

\textbf{Livrable 1 : Identification du Risque (Méthode EBIOS)}
\begin{itemize}
    \item Identifiez l'Actif Essentiel concerné.
    \item Identifiez la Menace (Source + Motivation).
    \item Identifiez la Vulnérabilité exploitée (citezen 2 failles : une humaine, une technique).
    \item Définissez l'Événement Redouté (Impact sur D.I.C.T).
\end{itemize}

\textbf{Livrable 2 : Estimation du Risque}
\begin{itemize}
    \item Proposez une Gravité (sur 5) et justifiez-la (Impact Business).
    \item Proposez une Probabilité (sur 5) et justifiez-la (Facilité).
    \item Calculez la Criticité et placez le risque sur une matrice.
\end{itemize}

\textbf{Livrable 3 : Plan d'Action (PAS)}
\begin{itemize}
    \item Proposez 3 mesures correctives immédiates pour réduire ce risque.
    \item Pour chaque mesure, indiquez l'option de traitement choisie (Réduction, Transfert, etc.).
    \item Définissez un RACI simple pour la mise en œuvre de ces mesures.
\end{itemize}

\textbf{Livrable 4 : Argumentaire Direction}
\begin{itemize}
    \item Rédigez un paragraphe de 5 lignes max à destination du PDG pour lui expliquer pourquoi ce risque, bien que l'attaque ait échoué, reste "Critique" et justifie des investissements immédiats.
\end{itemize}

% ==============================================================================
% Fichier : 03_module_audit.tex (Extrait Chapitre 1)
% Description : Introduction à l'Audit et au Management de la Qualité
% ==============================================================================